\section{The \sys{} Protocol}
\label{sec:protocol}

\sys{} gives command for a specified time, and a quorum of officers can extend that time.
The protocol has three primary properties. First, each device compares the end time of
command with its clock. Thus the officers remove a commander when they sign no more
renewals, and they send no message to the device. Second, a quorum of officers signs each
certificate, thus a small number of hostile officers has no effect. Third, a new commander
starts only after the end of each renewal that its signers know.

\subsection{Parameters and Certificates}

\Cref{tab:params} gives the parameters. Let $n \ge 3f + 1$ and
$q = \lceil (n + f + 1)/2 \rceil$. Then $f + 1$ or more officers are in each two sets of
$q$ officers, and one of these officers is loyal. Also $q \le n - f$, thus the loyal
officers can make a quorum. Last, $q > f$, thus the hostile officers cannot make a quorum
without loyal officers.

\begin{table}[t]
\caption{Parameters of \sys{}. The values are examples. They do not come from a power grid.}
\label{tab:params}
\centering
\small
\begin{tabularx}{\columnwidth}{@{}l>{\raggedright\arraybackslash}Xl@{}}
\toprule
Symbol & Name & Example \\
\midrule
$n$ & Number of officers & 4 \\
$f$ & Maximum number of hostile officers & 1 \\
$q$ & Quorum & 3 \\
$H$ & Horizon: the maximum time from a signature to the end time of a renewal & 600\,s \\
$g$ & Guard between two epochs, $g > 2\eps$ & 1\,s \\
$F$ & Fallback period of the deputy & 0\,s or 120\,s \\
$\eps$ & Clock error & 0.2\,s \\
$\Delta$ & Renewal interval & 30\,s \\
$\delta$ & Collection time for $q$ signatures & 2\,s \\
\bottomrule
\end{tabularx}
\end{table}

The protocol has two types of certificate. A \emph{grant} $\grant{k,h,b,s}$ gives command
in epoch $k$ to the commander $h$ from the start time $s$. It also gives the deputy $b$,
an officer other than $h$. A
\emph{renewal} $\renew{k,e}$ extends command in epoch $k$ to the end time $e$. A
certificate is \emph{valid} if it has the signatures of $q$ different officers on the same
contents. At the start, each party has the valid grant $\grant{1,h_1,b_1,0}$.

\subsection{Rules for an Officer}

Each loyal officer $u$ keeps two numbers. The first number is the epoch $E_u$, initially 1.
The second number is $M_u$, the largest end time in a renewal that $u$ signed, initially 0.
When $u$ receives a valid grant of an epoch $j > E_u$, it sets $E_u$ to $j$.

\begin{description}
  \item[O1] \emph{Renewal.} The officer $u$ signs $\renew{k,e}$ only if three
  conditions are correct. It has the valid grant of epoch $k$, and $k = E_u$. The end time is
  $e \le C_u + H$. The officer gives consent to the commander of epoch $k$. Then $u$ sets
  $M_u$ to the larger of $M_u$ and $e$.
  \item[O2] \emph{Grant.} The officer $u$ signs $\grant{j,h,b,s}$ only if three conditions are
  correct. The epoch is $j > E_u$, and $b \ne h$. The start time is $s \ge M_u + F + g$. The officer gives
  consent to $h$ as the commander. Then $u$ sets $E_u$ to $j$.
  \item[O3] \emph{Removal.} To remove the commander of epoch $k$, the officer $u$ stops its
  consent. Then $u$ signs no more renewals of epoch $k$.
\end{description}

Consent is a decision of the person. The protocol does not tell an officer when to stop
consent. For example, an officer can stop consent when the commander gives dangerous
orders. An officer can also stop consent when the officer cannot speak to the commander by
telephone.

\subsection{Rules for the Commander}

The commander $c$ of epoch $k$ sends a request for a renewal at intervals of $\Delta$. The
request has the end time $e = C_c + H - 2\eps$. Each officer who can sign by rule O1 signs
and sends its signature to the commander. When the commander has $q$ signatures, it has a
valid renewal. The commander attaches the grant and the newest valid renewal to each order.

\subsection{Rules for a Device}

A device keeps each valid certificate that it receives. Let $\tau$ be the clock time of the
device when it receives an order. The device accepts the order only if one of the rules D1
and D2 applies, and rule D3 does not stop it.

\begin{description}
  \item[D1] \emph{Commander.} The device accepts an order that $h$ signed for epoch $k$ if it has
  a valid $\grant{k,h,b,s}$ and a valid $\renew{k,e}$ with $s \le \tau \le e$.
  \item[D2] \emph{Deputy.} This rule applies only if $F > 0$. The device accepts an order that
  $b$ signed for epoch $k$ if three conditions are correct. It has a valid
  $\grant{k,h,b,s}$ with $s \le \tau$. It has one or more valid renewals of epoch $k$, and
  none of them has an end time $e \ge \tau$. The time is in the fallback interval:
  $a < \tau \le a + F$, where $a$ is the \emph{anchor} of epoch $k$.
  \item[D3] \emph{Local control.} The device has a switch with the positions ``REMOTE''
  and ``LOCAL''.
  Only a person at the device can set the switch. In the position ``LOCAL'', the device accepts
  no order from the network.
\end{description}

The device sets the anchor of an epoch one time. Before the device accepts the first order
of the deputy for epoch $k$, the anchor can change. It is the largest end time of the valid
renewals of epoch $k$ that the device has. When the device accepts that order, the anchor of epoch $k$
does not change again. Without this rule, an attacker can show new renewals to the device
one at a time. Then the deputy gets a new period $F$ after each renewal.

An attacker can replay an order of the commander while the certificates of that order are
valid. Thus a design must give a number to each order, and the device must discard an order
with a number that is not new. This is a known method, and the proofs in this paper do not
use it.

Rule D3 gives the input that \cref{lem:subst} does not include. A person at the device
can stop a hostile commander without a message from the officers. \pending{Not done at
this time: a power engineer must examine rule D3.}

\subsection{Change of the Commander}
\label{sec:change}

For a new commander, a grant with a higher epoch is necessary. An officer $p$ sends a proposal
$\grant{j,h',b',s}$ with $j = E_p + 1$ to the officers. The grant is valid when $q$
officers sign it. Three procedures use this step.

\paragraph{Removal}
The officers who want to remove the commander stop their consent (rule O3). When
$n - q + 1$ loyal officers do this, command expires at each device (\cref{thm:tenure}). A
new grant can use the start time $s = C_p + H + F + g + 2\eps$. Each loyal officer who
signed no renewal after the proposal can sign this grant by rule O2.

\paragraph{Prepared handover}
Let the officers prepare a handover for a time $T$. From the time $T - H$, the commander sends
each request with the end time $T$ or less. Then $M_u \le T$ for each officer $u$, and the
new grant can start at $T + F + g$. The officers sign the new grant after the last renewal,
because an officer who signs a grant signs no more renewals of the previous epoch (rule O1).
Thus a prepared handover has a period of $F + g$ between the two commanders, and not a
period of $H$.

\paragraph{Fallback to the deputy}
If $F > 0$, the deputy can give orders in a period $F$ after the last renewal expires at a
device (rule D2). The officers can use this period to make a new grant. If the officers do
not make a grant, the device accepts no order after this period.

\Cref{fig:timeline} shows the sequence at one device.

\begin{figure}[t]
\centering
\begin{tikzpicture}[x=1cm,y=1cm,font=\footnotesize,>=Stealth]
  % time axis
  \draw[->,ink] (0,0) -- (8.5,0) node[below left,inner sep=2pt] {time};
  % epoch k: commander
  \fill[accept!55] (0.4,0.25) rectangle (4.3,0.75);
  \node at (2.0,0.5) {commander $h$ (D1)};
  % deputy period
  \fill[pattern=north east lines,pattern color=fence] (4.3,0.25) rectangle (5.6,0.75);
  \draw[fence] (4.3,0.25) rectangle (5.6,0.75);
  \node[fill=white,inner sep=1pt] at (4.95,0.5) {$b$ (D2)};
  % next epoch
  \fill[accept!55] (6.2,0.25) rectangle (8.3,0.75);
  \node at (7.25,0.5) {commander $h'$};
  % renewals
  \foreach \x in {1.0,1.7,2.4} {\draw[ink] (\x,0.75) -- (\x,0.95);}
  \node[above] at (1.7,0.92) {renewals};
  % stop of consent
  \draw[->,warn,thick] (2.9,1.55) -- (2.9,0.8);
  \node[above,warn] at (2.9,1.5) {officers stop consent};
  % ticks
  \foreach \x/\l in {0.4/$s$,2.9/$t_0$,4.3/$e$,6.2/$s'$} {
    \draw[ink] (\x,0.08) -- (\x,-0.08);
    \node[below] at (\x,-0.08) {\l};
  }
  % braces
  \draw[decorate,decoration={brace,mirror,amplitude=3pt},ink] (2.9,-0.55) -- (4.25,-0.55)
    node[midway,below=3pt] {$\le H + 2\eps$};
  \draw[decorate,decoration={brace,mirror,amplitude=3pt},ink] (4.35,-0.55) -- (5.55,-0.55)
    node[midway,below=3pt] {$F$};
  \draw[decorate,decoration={brace,mirror,amplitude=3pt},ink] (5.65,-0.55) -- (6.2,-0.55)
    node[midway,below=3pt] {$g$};
  \node[above] at (7.25,0.78) {epoch $k{+}1$};
\end{tikzpicture}
\caption{Command at one device in epoch $k$ and in the subsequent epoch. From the time
$t_0$, the officers sign no more renewals. The commander $h$ keeps the device until the end
time $e$ of the last renewal. The deputy $b$ can give orders in the period $F$, if
$F > 0$. The new commander $h'$ starts after the guard $g$.}
\label{fig:timeline}
\end{figure}
